We conceptualize the design of RS as a contextual bandit problem where an agent interacts with a RS environment by taking an action (recommendation) based on a context (user) observed. More formally, the environment $\mathcal{E}$ is specified by a triple that consists of an observation space $\mathcal{O}$, an action space $\mathcal{A}$, and an observation probability function $P_O$. 

\textbf{Observation space $\mathcal{O}$}: In a contextual bandit problem, the observation generated by the environment offers feedback in response to the previous action taken by the agent as well as a new context the agent needs to execute an action. Formally, $\mathcal{O} = \mathcal{S} \times \mathbb{R}$ and at time step $t$, $O_t = (S_t, R_t)$. $S_t \in \mathcal{S}$ is a new context.  This context could, for example, provide a user's demographic information as well as data from their past interactions.  We will sometimes refer to $S_t$ as a user, since we think of it as providing information that distinguishes users.  $R_t \in \mathbb{R}$ denotes scalar feedback from the user of the previous time step $t-1$ in response to the action $A_{t-1}$.  For example, $R_t$ could be binary-valued, indicating whether the user clicked in response to a recommendation.  We will refer to $R_t$ as the reward received by the agent at time $t$.  

\textbf{Action space $\mathcal{A}$}: At each time $t$, the agent selects an action $A_{t} \in \mathcal{A}$, which identifies a content unit for recommendation to the user $S_t$. The action space $\mathcal{A}$ specifies the range of possibilities afforded by the universe of available content. 
% Given the action and observation definition, we denote History as $H_t = (A_0, O_1, A_1, O_2, A_2, \dots, O_t)$ for simplicity of notation in the rest of the paper. Note that $A_0$ is a dummy action which does not have any consequences from the environment. 

\textbf{Observation probability function $P_O$}: The observation probability function assigns a probability $P_O(O_{t+1}|S_t,A_t)$ to each possible outcome $O_{t+1} = (S_{t+1}, R_{t+1})$ from recommending $A_t$ to user $S_t$.  This probability is a product of two others, specified by functions $P_S$ and $P_R$:
\vspace{-3pt}
\begin{align}
    P_O(O_{t+1} | S_t, A_t) = P_S(S_{t+1})P_{R}(R_{t+1} | S_t, A_t).
\vspace{-5pt}
\end{align}
$P_S$ provides context probabilities, and $P_R$ provides reward probabilities.  At time $t$, the environment samples a new user $S_t \sim P_S$.  The agent then supplies a recommendation $A_t$, and the environment samples $R_{t+1} \sim P_{R}(\cdot| S_t, A_t)$.

The overall objective of the agent is to maximize its average reward $\frac{1}{T}\sum^T_{t=1} R_t$ over $T$ time steps. Note that, though each action $A_t$ does not directly impact any user other than $S_t$, the action influences what information is gathered through the observation $R_{t+1}$.  As such, the action bears an indirect impact on subequent users through what the agent learns, which can improve its ability to make useful recommendations. 
 Optimizing the delayed benefits of learning requires strategic sequential decisions that balance exploration versus exploitation.

We think of $S_t$ and $A_t$ as offering information about users and content units in arbitrary raw formats.  Agents, however, often rely on starting with a more structured representations provided by domain experts.  In this work, we also assume availability of feature extractors that map $S_t$ and $A_t$ to such representations $\phi_{A_t}$ and $\psi_{S_t}$.

Note that our problem formulation encompasses only online learning.  In a real RS, one would pretrain an agent offline on historical data before engaging in online learning.  The methods we develop for online learning can be applied in this workflow post pretraining of an epinet-enhanced model.  However, in order to focus our attention on the problem of exploration, in this paper, we limit our discussions, designs and experiments to the online learning part and leave offline learning for future work.